\documentclass[11pt]{article}
\usepackage[margin=1in]{geometry}
\usepackage{amsmath,amssymb,amsthm}
\usepackage[hidelinks]{hyperref}

\newtheorem{theorem}{Theorem}
\newtheorem{proposition}[theorem]{Proposition}
\newtheorem{lemma}[theorem]{Lemma}
\newtheorem{corollary}[theorem]{Corollary}
\theoremstyle{definition}
\newtheorem{remark}[theorem]{Remark}
\newtheorem{definition}[theorem]{Definition}

\newcommand{\R}{\mathbb{R}}
\newcommand{\norm}[1]{\left\|#1\right\|}
\newcommand{\abs}[1]{\left|#1\right|}
\renewcommand{\d}{\,d}

\title{Route A: full quantitative proof of sharp Faber--Krahn stability,\\
audit-ready end-to-end version}
\author{Plan 1 / Sharp stability for Faber--Krahn}
\date{2026-05-11 (final pass)}

\begin{document}
\maketitle

\begin{abstract}
This note assembles the full end-to-end quantitative proof of sharp
Faber--Krahn stability. Plan 1's stated goal is to make BDV's proof
quantitative by removing the contradiction step. The single-set selection
map of \texttt{quantitative-selection-principle.md} is the device that
achieves this. The remainder of the chain (graph entry, BDV nearly
spherical, Kohler--Jobin) is constructive and supplies explicit constants.

Every lemma is stated with explicit hypotheses, every constant is named with
its $(N,R)$-dependence noted, and every reference to BDV's free-boundary
paper is given by theorem number. The argument as a whole has no
contradiction step.

\medskip
\noindent
\textbf{Scope:} this note documents Route A (using BDV Theorem~3.3 as the
final closure). Route B (replacing BDV Theorem~3.3 by the Bernoulli
spectral closure of \texttt{fixed-domain-bernoulli-expansion.md}) is an
optional alternative and is not used here. Wherever a step uses a published
BDV theorem as a black box, it is stated as such.
\end{abstract}

\tableofcontents

\section{Notation and standing hypotheses}\label{sec:notation}

Throughout, $N\ge2$, $0<\gamma<1$ are fixed. $B_R$ denotes the open ball
of radius $R$ centered at the origin in $\R^N$, $\abs{B_R}=\omega_NR^N$,
$\omega_N=\abs{B_1}$. All sets $\Omega$ considered are open with finite
Lebesgue measure.

\subsection{Functionals}

\paragraph{Torsion potential and energy.}
For an open bounded $\Omega$, the torsion potential $u_\Omega$ is the unique
weak solution of
\[
-\Delta u_\Omega=1\hbox{ in }\Omega,\qquad u_\Omega\in H^1_0(\Omega).
\]
The torsional rigidity is $T(\Omega)=\int_\Omega u_\Omega\d x
=\int_\Omega\abs{\nabla u_\Omega}^2\d x$. The torsion energy is
\[
E(\Omega):=\tfrac12\int_\Omega\abs{\nabla u_\Omega}^2-\int_\Omega u_\Omega
=-\tfrac12 T(\Omega).
\]
For comparison ball $B$ at fixed volume,
$E(\Omega)\ge E(B)$ (Saint--Venant), with equality iff $\Omega=B$
(modulo null sets and translations).

\paragraph{Faber--Krahn data.}
$\lambda_1(\Omega)$ is the first Dirichlet eigenvalue of $-\Delta$ on
$\Omega$; Faber--Krahn: $\lambda_1(\Omega)\ge\lambda_1(B^*)$ at fixed
volume, equality iff ball.

\paragraph{Asymmetry.}
The Fraenkel asymmetry is
\[
\mathcal A(\Omega):=
\inf_{x_0\in\R^N}\frac{\abs{\Omega\Delta(x_0+B^*)}}{\abs{\Omega}},
\]
$B^*$ the ball of volume $\abs{\Omega}$ at the origin.

\paragraph{Quadratic asymmetry.}
The integrated radial asymmetry used by BDV is
\[
\alpha(\Omega):=
\int_{\Omega\Delta B_{r_\Omega}(x_\Omega)}
\bigl|r_\Omega-\abs{x-x_\Omega}\bigr|\,\d x,
\qquad
r_\Omega:=(\abs{\Omega}/\omega_N)^{1/N},
\]
where $x_\Omega$ is the barycenter. For unit volume this is BDV Definition~4.1.
BDV Lemma~4.2 gives $\mathcal A(\Omega)^2\le C_{\rm Fr}(N)\alpha(\Omega)$.

\paragraph{Quotient.}
\[
Q_\alpha(\Omega):=\frac{E(\Omega)-E(B_1)}{\alpha(\Omega)}.
\]
The whole work is about lower bounds on $Q_\alpha$ in the small-asymmetry regime.

\subsection{BDV regularity package}

We use the BDV \emph{regularity package} as a black box: it asserts that the
minimizers of the penalized functionals below satisfy two-sided density,
Lipschitz, nondegeneracy, and the Bernoulli coefficient bounds needed to start
the flatness theory, with constants depending only on $(N,R)$. After graph
entry, the explicit selected Bernoulli law is smooth on the fixed collar and
the same constants control the Schauder bootstrap. Specifically:

\begin{itemize}
\item \textbf{BDV Lemma 4.9} (density). There are
$c_d(N,R)>0$, $r_d(N,R)>0$ such that for every penalized minimizer $U$ and
every free boundary point $x\in\partial U$ and $r\le r_d$,
\[
c_d r^N\le\abs{U\cap B_r(x)}\le(1-c_d)r^N\cdot\omega_N.
\]

\item \textbf{BDV Lemma 4.12} (Lipschitz/nondegeneracy). The torsion
function $u$ of a penalized minimizer is Lipschitz with constant $L_u(N,R)$
and nondegenerate: $\sup_{B_r(x)}u\ge c_d r$ for $x\in\partial U$, $r\le r_d$.

\item \textbf{BDV Lemma 4.16} (Bernoulli coefficient). For a penalized
minimizer in the flat regime, $q_u(y)=\abs{\nabla u(y)}$ is well-defined on
$\partial U$ and satisfies
\[
q_-(N,R)\le q_u\le q_+(N,R),\qquad[q_u]_{C^{1,\gamma}(\partial U)}\le L_q(N,R).
\]

\item \textbf{BDV Theorem 4.18} (Alt--Caffarelli flatness improvement). If
the free boundary of a penalized minimizer is $\bar\mu(N,R)$-flat at scale
$\bar\rho(N,R)$, then it improves to a $C^{1,\gamma}$ graph at scale
$\bar\rho/2$ with norm $C_{\rm AC}(N,R)\bar\mu$.
\end{itemize}

These four results are quoted from BDV~\S 4. We use each one \emph{per single
set}, not along a sequence, and this is legitimate at the level of BDV's own
proofs: the proofs of Lemmas~4.9--4.12 open with ``being $j$ fixed we drop the
subscript'' and proceed by a one-set variational comparison against the
perturbed minimality inequality~(4.19), which is itself the single-set first
variation of Lemma~4.2; the resulting constants depend only on $(N,R)$ and on
the regularity regime $q\le q_{\rm sel}(N,R)$ (equivalently BDV's
$\sigma\le\sigma_{\rm reg}(N,R)$, the penalization prefactor being
$\le\tau=q^{1/4}$). The \emph{only} place BDV pass to $j\to\infty$ is
Lemma~4.13 (Kuratowski convergence $\partial\Omega_j\to\partial B_1$), invoked
solely to produce a free-boundary point near every $\bar x\in\partial B_1$.
That single qualitative step is exactly what Lemma~\ref{lem:alphaH} below
replaces quantitatively, so Lemma~4.13 is \emph{not invoked} anywhere in this
route. The single-set quasi-minimality these per-set lemmas require is the one
Lemma~\ref{lem:rescale} secures for the selected $U_*$ before dilation.

\paragraph{BDV nearly spherical theorem (used as black box).}
The closing step is:
\begin{quote}
\textbf{BDV Theorem 3.3} (nearly spherical). For every $0<\gamma_0\le1$
there exist $c_{\rm sph}(N)>0$ and $\delta_{\rm sph}(N,\gamma_0)>0$ such that for every set
$U=\{(1+g(\theta))\theta:\theta\in\partial B_1\}$ with
$\abs{U}=\abs{B_1}$, $\hbox{barycenter}(U)=0$, and
$\norm{g}_{C^{2,\gamma_0}(\partial B_1)}\le\delta_{\rm sph}$,
\[
E(U)-E(B_1)\ge c_{\rm sph}(N)\norm{g}_{H^{1/2}(\partial B_1)}^2.
\]
\end{quote}
This is the only step that Route A uses BDV for in the final closure; Route B
would replace it by \texttt{fixed-domain-bernoulli-expansion.md}, Theorem 6.2.

\section{The BDV contradiction step and its single-set replacement}

\subsection{BDV's original contradiction}

BDV's proof (BDV~\S 5) of sharp Saint--Venant runs:
\begin{enumerate}
\item Assume by contradiction the sharp inequality fails: there is a
sequence $\{\Omega_n\}_{n\ge1}$ in $B_R$ with $\abs{\Omega_n}=\abs{B_1}$ and
\(Q_\alpha(\Omega_n)\to 0\).
\item For each $n$, apply a sequence of penalized minimizations along the
sequence to extract a \emph{selected sequence} $\{U_n\}$, with quantitative
control on the energy and asymmetry losses at each step.
\item Pass to a Hausdorff limit $U_\infty$ via the BDV regularity package
(equicontinuity of free boundaries). $U_\infty$ is a penalized minimizer
with $\alpha(U_\infty)=\lim\alpha(U_n)>0$.
\item Apply BDV Theorem~4.18 + nearly spherical to conclude $U_\infty=B_1$
(after barycenter), contradicting $\alpha(U_\infty)>0$.
\end{enumerate}

Step~3 is the contradiction step: the limit $U_\infty$ is qualitative and
the constants underlying it are not exposed.

\subsection{Single-set replacement: the selection map of
\texttt{quantitative-selection-principle.md}}

\begin{definition}[Single-set penalized functional]\label{def:penal}
For $\Omega_0\subset B_R$, $\abs{\Omega_0}=\abs{B_1}$, set
$\varepsilon=\alpha(\Omega_0)$, $q=Q_\alpha(\Omega_0)$. For
$\tau>0$ and $\hat\eta>0$, the single-set penalized functional is
\[
J_{\tau,\hat\eta,\Omega_0}(U)
:=E(U)+f_{\hat\eta}(\abs U)
+\sqrt{\varepsilon^2+\tau^2(\alpha(U)-\varepsilon)^2},
\]
$U\subset B_R$, where $f_{\hat\eta}(s)$ is the BDV symmetric volume
penalization
\[
f_{\hat\eta}(s):=
\begin{cases}
\hat\eta(s-\omega_N),&s\le\omega_N,\\
(s-\omega_N)/\hat\eta,&s\ge\omega_N,
\end{cases}
\]
with $\hat\eta=\hat\eta(N,R)$ chosen as in BDV Lemma~4.5. This penalty does
not force exact unit volume; the selected set is volume-normalized in
\S\ref{sec:volnorm}.
\end{definition}

\begin{theorem}[Single-set selection, \texttt{quantitative-selection-principle.md} Thm.~1]\label{thm:selection}
There exist $\tau_{\rm ex}(N,R)>0$, $\tau_{\rm reg}(N,R)>0$, and
$C_{\rm sel}(N,R)>0$ such that for every $\tau\le\tau_{\rm reg}(N,R)$,
$q\le\min(q_{\rm sel}(N,R),\tau^2)$, and every $\Omega_0\subset B_R$ with
$\abs{\Omega_0}=\abs{B_1}$ and $Q_\alpha(\Omega_0)=q$, the functional
$J_{\tau,\hat\eta,\Omega_0}$ admits a minimizer $U_*$ with the following
properties. After volume normalization (\S\ref{sec:volnorm} below) to
$\widetilde\Omega=r_*^{-1}U_*$, $r_*=(\abs{U_*}/\abs{B_1})^{1/N}$:
\begin{enumerate}
\item \emph{Existence and regularity.} $U_*$ is a BDV penalized minimizer
satisfying Lemmas~4.9, 4.12, 4.16 with constants depending only on
$(N,R)$. In particular, $\sigma_*:=$ penalization parameter of $U_*$
satisfies
\(
\sigma_*\le C\tau.
\)
\item \emph{Asymmetry preservation.} With
$\rho(q,\tau)=\sqrt{2q+q^2}/\tau+C_{\rm sel}q$,
\[
(1-\rho(q,\tau))\,\alpha(\Omega_0)
\le \alpha(\widetilde\Omega)
\le (1+\rho(q,\tau))\,\alpha(\Omega_0).
\]
\item \emph{Quotient preservation.}
\[
Q_\alpha(\widetilde\Omega)\le\frac{C_{\rm sel}(N,R)}{1-\rho(q,\tau)}\,Q_\alpha(\Omega_0).
\]
\end{enumerate}
With the canonical choice $\tau=q^{1/4}$, $\rho(q,\tau)\le1/2$ for
$q\le q_{\rm sel}$.
\end{theorem}

\begin{remark}
Theorem~\ref{thm:selection} is the \emph{exact single-set analogue} of BDV's
sequential selection. The only change is that BDV minimize a sequence of
functionals along a contradiction sequence, while here we minimize one
functional given one input set. The functional uses $\varepsilon=\alpha(\Omega_0)$
as a fixed parameter, not as a sequence.
\end{remark}

\begin{proof}[Proof sketch (full proof in \texttt{quantitative-selection-principle.md}).]
Existence: this is BDV Lemma~4.6 with the sequence parameter replaced by the
fixed parameter $\tau$. The proof is performed in the quasi-open class:
one truncates a minimizing sequence of torsion potentials, obtains a uniform
perimeter bound by the coarea argument, and then passes to the $L^1$/quasi-open
limit. The density estimates are consequences of minimality and are not used
to prove existence.

Regularity: the smooth Bernoulli term
$\sqrt{\varepsilon^2+\tau^2(\alpha(U)-\varepsilon)^2}$
contributes a $C^{0,1}$ pseudonorm bounded by $\tau$, and so falls within
the BDV regularity framework with selection parameter $\sigma\le C\tau$.

Quotient preservation: by minimization, $J_{\tau,\hat\eta,\Omega_0}(U_*)\le
J_{\tau,\hat\eta,\Omega_0}(\Omega_0)$. The volume penalization $f_{\hat\eta}$
forces $\abs{U_*-B_1}\le O(\delta_T)$, $\delta_T=E(U_*)-E(B_1)$ (BDV trick).
The $\alpha$-bracket gives
$\alpha(U_*)\ge\varepsilon-\rho\varepsilon$ with $\rho$ as in (ii), and the
quotient bound (iii) follows by combining (ii) with the energy bound
$E(U_*)\le E(\Omega_0)+\hbox{error}\le E(B_1)+(1+O(\tau))(E(\Omega_0)-E(B_1))$.
\end{proof}

\section{Volume normalization}\label{sec:volnorm}

The selected $U_*$ has volume $\abs{B_1}+O(\delta_T)$, so it must be
rescaled to compare with the unit ball.

\begin{lemma}[Volume-normalized law and quasi-minimality]\label{lem:rescale}
Set $r_*=(\abs{U_*}/\abs{B_1})^{1/N}$, $\widetilde\Omega=r_*^{-1}U_*$,
$\widetilde u(x)=r_*^{-2}u_{U_*}(r_*x)$. Then
\[
\abs{r_*-1}\le C\delta_T,\qquad
\abs{\widetilde\Omega}=\abs{B_1},\qquad
-\Delta\widetilde u=1\hbox{ in }\widetilde\Omega,\;\widetilde u=0\hbox{ on }\partial\widetilde\Omega,
\]
and the rescaled Bernoulli law on $\partial\widetilde\Omega$ reads
\[
\abs{\nabla\widetilde u(y)}^2=\widetilde\Lambda+\widetilde a_\sigma\widetilde\Psi(y),
\]
with $\widetilde\Lambda=r_*^{-2}\Lambda$, $\widetilde a_\sigma=r_*^{-1}a_\sigma$,
and $\widetilde\Psi$ the rescaled $\alpha$-bracket. Moreover, if $U_*$
satisfies BDV's perturbed minimality inequality with parameter $\sigma$, then
$\widetilde u$ satisfies the same type of inequality for the rescaled volume
penalty
\[
\widetilde f(s):=r_*^{-(N+2)}f_{\hat\eta}(r_*^Ns)
\]
and with quasi-minimality parameter $C\sigma r_*^{-2}$. In the regime
$\delta_T\le c$ small, $r_*\in[1/2,2]$, so all constants degrade by bounded
$(N,R)$-factors.
\end{lemma}

\begin{proof}
The torsion equation scales by $u\mapsto r_*^{-2}u(r_*\cdot)$ and the gradient
by $\nabla\widetilde u(y)=r_*^{-1}\nabla u_{U_*}(r_*y)$.

Write the first-variation bracket in the translation-invariant averaged form
\[
A_U:=\frac1{\abs U}\int_U\frac{z-x_U}{\abs{z-x_U}}\d z,\qquad
\Psi_U(x):=\abs{x-x_U}-A_U\cdot(x-x_U).
\]
This is BDV's bracket, up to an additive constant if one writes $A_U\cdot x$;
that constant is absorbed into $\Lambda$. Under the change of variables
$U_*=r_*\widetilde\Omega$ one has $A_{U_*}=A_{\widetilde\Omega}$ and
\[
\Psi_{U_*}(r_*y)=r_*\Psi_{\widetilde\Omega}(y).
\]
Thus
\[
\abs{\nabla\widetilde u(y)}^2
=r_*^{-2}\bigl(\Lambda+a_\sigma\Psi_{U_*}(r_*y)\bigr)
=r_*^{-2}\Lambda+r_*^{-1}a_\sigma\Psi_{\widetilde\Omega}(y).
\]

For the quasi-minimality, let $\zeta$ be an admissible competitor for
$\widetilde u$ and set $v(x)=r_*^2\zeta(x/r_*)$. Then
\[
\mathcal E(v)=r_*^{N+2}\mathcal E(\zeta),\qquad
\abs{\{v>0\}}=r_*^N\abs{\{\zeta>0\}},
\]
and
\[
\abs{U_*\Delta\{v>0\}}=r_*^N
\abs{\widetilde\Omega\Delta\{\zeta>0\}}.
\]
Dividing BDV's perturbed minimality inequality for $U_*$ by $r_*^{N+2}$
gives the displayed rescaled inequality.
\end{proof}

\begin{remark}[Proof order]
The free-boundary regularity arguments below are applied to the original
selected minimizer $U_*$ before the final dilation. Lemma~\ref{lem:rescale}
shows that applying them after dilation would also be legitimate, but the
pre-dilation order avoids any hidden use of post-normalization minimality.
\end{remark}

\section{Graph entry}\label{sec:graphentry}

The selected minimizer is regularized before the final dilation; equivalently,
by Lemma~\ref{lem:rescale}, the volume-normalized set
$\widetilde\Omega$ is a BDV-type quasi-minimizer with uniformly transported
constants. It becomes a global graph over $\partial B_1$ once its asymmetry is
below an explicit threshold.

\subsection{Asymmetry to Hausdorff distance}

\begin{lemma}[Asymmetry to Hausdorff distance]\label{lem:alphaH}
For every volume-normalized selected quasi-minimizer $\widetilde\Omega$ with
$\abs{\widetilde\Omega}=\abs{B_1}$ and comparison ball $B_1$ (after
translating by barycenter),
\[
d_H(\partial\widetilde\Omega,\partial B_1)
\le C_H'(N,R)\,\alpha(\widetilde\Omega)^{1/(2N)}.
\]
\end{lemma}

\begin{proof}
Write $\alpha:=\alpha(\widetilde\Omega)$ and
$d:=d_H(\partial\widetilde\Omega,\partial B_1)$. Two per-set ingredients enter
-- the asymmetry--volume comparison and the two-sided density estimate --
and no limit is taken anywhere.

\emph{Step 1 (asymmetry controls the symmetric difference).} By BDV Lemma~4.2,
for a penalized minimizer normalized by its barycenter with
$\abs{\widetilde\Omega}=\abs{B_1}$,
\[
\abs{\widetilde\Omega\Delta B_1}\le C_\alpha(N)\,\alpha^{1/2},
\]
the exponent $1/2$ accounting for the passage from the optimal Fraenkel
center to the barycenter-centered ball.

\emph{Step 2 (density turns volume into one-sided distance).} Fix
$x\in\partial\widetilde\Omega$ and set $t:=\operatorname{dist}(x,\partial B_1)$.
Since $B_t(x)$ does not meet $\partial B_1$ it lies either inside $B_1$ or in
$\R^N\setminus\overline{B_1}$. Put $r:=\min\{t,r_d\}$ with $r_d=r_d(N,R)$ the
density radius of the package. If $B_t(x)\subset B_1$, the lower density of the
\emph{complement} at the free-boundary point $x$ (the upper-density half of
BDV Lemma~4.9) gives
\[
\abs{B_1\setminus\widetilde\Omega}
\ \ge\ \abs{\widetilde\Omega^{\,c}\cap B_r(x)}
\ \ge\ c_d\,\omega_N r^{N};
\]
if instead $B_t(x)\subset\R^N\setminus\overline{B_1}$, the lower density of
$\widetilde\Omega$ gives $\abs{\widetilde\Omega\setminus B_1}\ge
c_d\,\omega_N r^{N}$. In either case
\[
\abs{\widetilde\Omega\Delta B_1}\ \ge\ c_d\,\omega_N\,
\bigl(\min\{t,r_d\}\bigr)^{N}.
\]
Once $\alpha\le\alpha_0(N,R)$ is small enough that
$C_\alpha\alpha_0^{1/2}<c_d\,\omega_N r_d^{\,N}$, the minimum cannot be
attained at $r_d$, so it is attained at $t$ and, with Step~1,
\[
\sup_{x\in\partial\widetilde\Omega}\operatorname{dist}(x,\partial B_1)
\ \le\ \Bigl(\tfrac{C_\alpha}{c_d\,\omega_N}\Bigr)^{1/N}\alpha^{1/(2N)} .
\]

\emph{Step 3 (the reverse inclusion, no density needed).} Let
$\bar x\in\partial B_1$ and $s:=\operatorname{dist}(\bar x,\partial\widetilde
\Omega)$. The ball $B_s(\bar x)$ contains no point of
$\partial\widetilde\Omega$, hence lies entirely in $\widetilde\Omega$ or
entirely in its complement. For $s\le\tfrac12$ each of the two caps
$B_s(\bar x)\cap B_1$ and $B_s(\bar x)\setminus\overline{B_1}$ has volume at
least $c_N s^{N}$, $c_N=c_N(N)>0$. If $B_s(\bar x)\subset\widetilde\Omega$ the
outer cap lies in $\widetilde\Omega\setminus B_1$; if
$B_s(\bar x)\subset\widetilde\Omega^{\,c}$ the inner cap lies in
$B_1\setminus\widetilde\Omega$. Either way $\abs{\widetilde\Omega\Delta
B_1}\ge c_N s^{N}$, so $s\le(C_\alpha/c_N)^{1/N}\alpha^{1/(2N)}$.

Combining Steps~2--3,
\[
d_H(\partial\widetilde\Omega,\partial B_1)\le C_H'(N,R)\,\alpha^{1/(2N)},
\qquad
C_H'=\max\Bigl\{\bigl(\tfrac{C_\alpha}{c_d\omega_N}\bigr)^{1/N},
\bigl(\tfrac{C_\alpha}{c_N}\bigr)^{1/N}\Bigr\},
\]
for every $\alpha\le\alpha_0(N,R)$, where $\alpha_0$ is the minimum of the
finitely many displayed smallness conditions
($C_\alpha\alpha_0^{1/2}<c_d\omega_N r_d^{\,N}$ for Step~2, and
$(C_\alpha/c_N)^{1/N}\alpha_0^{1/(2N)}\le\tfrac12$ so that $s\le\tfrac12$ in
Step~3). The threshold $\alpha_0$ is thus a fixed $(N,R)$-quantity (it is
absorbed into $\alpha_{\rm graph}$ in \S\ref{sec:graphentry}); it depends on
neither $\delta_T$ nor the conclusion.
The constants $c_d,r_d$ are the per-set density constants of the regularity
package, valid for the single selected set by the discussion following its
statement; no sequence or compactness is used.
\end{proof}

\subsection{Hausdorff to flatness}

Pick BDV Theorem~4.18 thresholds $\bar\mu(N,R),\bar\rho(N,R)$. Define
$\mu(N,R)\le\bar\mu/16C_0$ small enough for the rest of the chain
(\S\ref{sec:bootstrap}), where $C_0$ is the universal conversion constant
between the Alt--Caffarelli slab and the $F(\cdot,1,+\infty)$ normalization.
Pick $\rho_\mu\le\bar\rho$ so that every cap of $\partial B_1$ of radius
$6\rho_\mu$ lies in a $\mu\rho_\mu$-slab around its tangent. Set
$h_F:=\mu\rho_\mu/16$.

\begin{proposition}[Hausdorff closeness to Alt--Caffarelli flatness]\label{prop:flatness}
If $d_H(\partial\widetilde\Omega,\partial B_1)\le h_F$, then for every
$\bar x\in\partial B_1$ there is $x_0\in\partial\widetilde\Omega\cap B_{h_F}(\bar x)$
such that in $B_{4\rho_\mu}(x_0)$, the torsion function $\widetilde u$ is of
Alt--Caffarelli class $F(C_0\mu,1,+\infty)$ with respect to the tangent
direction $\nu_{\bar x}$.
\end{proposition}

\begin{proof}
First, Hausdorff closeness and $\abs{\widetilde\Omega}=\abs{B_1}$ give the
annular sandwich
\[
B_{1-h_F}\subset\widetilde\Omega\subset B_{1+h_F}.
\]
The outer inclusion follows because a point of $\widetilde\Omega$ outside
$B_{1+h_F}$ would force a boundary point farther than $h_F$ from $\partial B_1$.
For the inner inclusion, any component of
$B_{1-h_F}\setminus\widetilde\Omega$ would have boundary inside $B_{1-h_F}$,
again contradicting the Hausdorff bound; the alternative
$B_{1-h_F}\subset\R^N\setminus\widetilde\Omega$ contradicts the volume identity
for $h_F$ small, since $\widetilde\Omega\subset B_{1+h_F}$.

Choose $x_0\in\partial\widetilde\Omega\cap B_{h_F}(\bar x)$ by the
Hausdorff bound. The curvature of $\partial B_1$ gives
\[
\partial B_1\cap B_{5\rho_\mu}(\bar x)\subset
\{x:\abs{(x-\bar x)\cdot\nu_{\bar x}}\le C\rho_\mu^2\}.
\]
Combining this with the Hausdorff bound and recentering at $x_0$ yields
\[
\partial\widetilde\Omega\cap B_{4\rho_\mu}(x_0)
\subset
\{x:\abs{(x-x_0)\cdot\nu_{\bar x}}\le C(\rho_\mu^2+h_F)\}.
\]
The choices $\rho_\mu\le c\mu$ and $h_F\le c\mu\rho_\mu$ make the last height
at most $C\mu\rho_\mu$.

With $\nu_{\bar x}$ oriented as the inward normal to $B_1$, the slab bound and
the annular sandwich pin \emph{both} phases deterministically. A point
$x\in B_{4\rho_\mu}(x_0)$ with
$(x-x_0)\cdot\nu_{\bar x}\le -C_0\mu(4\rho_\mu)$ lies outside $B_{1+h_F}$ once
$C_0$ is chosen large, hence, by the outer inclusion
$\widetilde\Omega\subset B_{1+h_F}$ of the sandwich, satisfies
$\widetilde u(x)=0$. Symmetrically, a point with
$(x-x_0)\cdot\nu_{\bar x}\ge +C_0\mu(4\rho_\mu)$ lies inside $B_{1-h_F}$, hence,
by the inner inclusion $B_{1-h_F}\subset\widetilde\Omega$, lies in
$\widetilde\Omega$, so $\widetilde u>0$ there. Thus the positive phase of the
class $F(C_0\mu,1,+\infty)$ is \emph{not} an additional hypothesis to be
assumed: it is supplied, with the same explicit constants, by the inner
inclusion of the sandwich, exactly mirroring the empty outer half-slab.

The only quantitative input the flatness-improvement iteration consumes on the
positive side is a growth rate, and this is the \emph{explicit per-set}
nondegeneracy of the regularity package: for every $x_0\in\partial\widetilde
\Omega$ and $r\le r_d(N,R)$,
\[
\sup_{B_r(x_0)}\widetilde u\ \ge\ c_d(N,R)\,r ,
\]
(BDV Lemma~4.12; the constant is $c_d(N,R)$, never relative to $\delta_T$),
together with the two-sided density $c_d\,\omega_N r^{N}\le
\abs{\widetilde\Omega\cap B_r(x_0)}\le(1-c_d)\,\omega_N r^{N}$ of BDV
Lemma~4.9. Both hold for the single selected set under
$q\le q_{\rm sel}(N,R)$, by the per-set discussion in
\S\ref{sec:notation}. With the zero set, the positive set, and the
nondegeneracy/density constants all exhibited, $\widetilde u$ is of
Alt--Caffarelli class $F(C_0\mu,1,+\infty)$ in $B_{4\rho_\mu}(x_0)$ with the
package's $(N,R)$-constants.
\end{proof}

\subsection{Local flatness to global graph}

\begin{proposition}[Local flatness to a global radial graph]\label{prop:globalgraph}
Assume Prop.~\ref{prop:flatness} holds at every $\bar x\in\partial B_1$.
Choose $\mu$ small enough that $C_0\mu\le\bar\mu$ (so BDV~4.18 applies) and
$C_{\rm AC}\mu\le\mu_{\rm tub}(N)$, the tubular-neighborhood slope for radial
projection. Then
\[
\partial\widetilde\Omega=\{(1+g(\theta))\theta:\theta\in\partial B_1\}
\]
with
\[
\norm{g}_{L^\infty(\partial B_1)}\le C(N,R)\,d_H(\partial\widetilde\Omega,\partial B_1),
\quad
\norm{g}_{C^{1,\gamma}(\partial B_1)}\le C(N,R)\mu.
\]
\end{proposition}

\begin{proof}
BDV Theorem 4.18 in each $B_{\rho_\mu}(x_0)$ of Prop.~\ref{prop:flatness}
gives local $C^{1,\gamma}$ graphs of norm $\le C_{\rm AC}\mu$. These charts
cover the entire free boundary: if $y\in\partial\widetilde\Omega$, choose
$\bar x=y/\abs y\in\partial B_1$; the annular sandwich gives
$\abs{y-\bar x}\le h_F$, while the corresponding chart center $x_0$ satisfies
$\abs{x_0-\bar x}\le h_F$, so $y\in B_{\rho_\mu}(x_0)$ after shrinking
$h_F/\rho_\mu$.

On every local chart the radial projection is a $C^{1,\gamma}$ diffeomorphism
onto its image, because the local normal makes an angle bounded away from
$\pi/2$ with the radial direction when $C_{\rm AC}\mu\le\mu_{\rm tub}$.
Surjectivity of the global radial projection follows from the annular
sandwich: every ray from the origin meets $\partial\widetilde\Omega$ between
radii $1-h_F$ and $1+h_F$. For injectivity, if two boundary points lay on the
same ray, their distance would be at most $2h_F<\rho_\mu/2$, so both would lie
in one local chart; a transverse local graph is intersected by such a radial
line at most once. Hence the intersection on every ray is unique.

The unique radial intersection defines $g$. The $L^\infty$ bound follows from
the annular sandwich, and the $C^{1,\gamma}$ bound follows by passing the
finite family of local chart estimates through the uniformly nondegenerate
radial-projection coordinates.
\end{proof}

\subsection{Graph-entry threshold}

\begin{theorem}[Graph-entry threshold]\label{thm:graphentry}
Set
\[
\alpha_{\rm graph}(N,R):=\Bigl(\frac{h_F}{C_H'(N,R)}\Bigr)^{2N}.
\]
Every volume-normalized selected quasi-minimizer $\widetilde\Omega\subset B_R$ with
$\abs{\widetilde\Omega}=\abs{B_1}$, barycenter $0$, and
\(
\alpha(\widetilde\Omega)\le\alpha_{\rm graph}
\)
is a $C^{1,\gamma}$ normal graph over $\partial B_1$ with
$\norm{g}_{C^{1,\gamma}}\le C(N,R)\mu(N,R)$.
\end{theorem}

\begin{proof}
Combine Lemmas~\ref{lem:alphaH}, Propositions~\ref{prop:flatness},
\ref{prop:globalgraph}.
\end{proof}

\section{Uniform Schauder bounds and small $C^{2,\gamma_0}$ entry}\label{sec:bootstrap}

BDV Theorem~3.3 needs the graph norm to be small in $C^{2,\gamma_0}$ for some
$0<\gamma_0\le1$. Graph entry gives small $C^{1,\gamma}$ norm and, more
importantly, small $L^\infty$ norm through Hausdorff closeness. The Schauder
bootstrap supplies uniform higher norms; interpolation then turns small
$L^\infty$ plus uniform high regularity into small $C^{2,\gamma_0}$.

\begin{proposition}[Uniform bootstrap]\label{prop:bootstrap}
Let $\partial\widetilde\Omega=\{(1+g)\theta\}$ be a graph obtained from
Theorem~\ref{thm:graphentry}, and let $\widetilde u$ be its torsion potential.
For every integer $m\ge3$ there is $M_m(N,R)<\infty$ such that
\[
\norm{g}_{C^{m,\gamma}(\partial B_1)}\le M_m(N,R).
\]
\end{proposition}

\begin{proof}
Work first with the unscaled selected minimizer $U_*$, and write
$\abs{U_*}=\abs{B_r}$, with $r\in[1/2,2]$ in the small-deficit regime. The
scale-invariant form of the graph-entry step gives a $C^{1,\gamma}$ graph over
$\partial B_r(x_{U_*})$ and
the annular containment
\[
B_{r-h}(x_{U_*})\subset U_*\subset B_{r+h}(x_{U_*}),\qquad h\le r/4.
\]
Hence a fixed collar of the free boundary lies in
$\{r/2\le\abs{x-x_{U_*}}\le3r/2\}$.

The selected Bernoulli coefficient has the form
\[
q_{U_*}(x)^2=\Lambda+a_\sigma\Psi_{U_*}(x),\qquad
\Psi_{U_*}(x)=\abs{x-x_{U_*}}-A_{U_*}\cdot(x-x_{U_*}),
\]
with $\abs{A_{U_*}}\le1$, $\abs{a_\sigma}\le C\sigma$, and
$0<q_-\le q_{U_*}\le q_+$. On the fixed annular collar,
\[
\norm{\Psi_{U_*}}_{C^k}\le C_k(N,R)
\]
for every fixed $k$, since only derivatives of $\abs{x-x_{U_*}}$ and a linear
term occur. Choosing $\sigma$ below BDV's regularity threshold and using the
positive lower bound on $q_{U_*}$ gives
\[
\norm{q_{U_*}}_{C^k(\hbox{\rm collar})}\le C_k(N,R).
\]
This is the single-set content of BDV Lemma~4.16; no derivative of the unknown
graph enters this estimate.

Now apply the higher-regularity clause of BDV Theorem~4.18. The local
Alt--Caffarelli charts obtained in Theorem~\ref{thm:graphentry} have
$C^{m,\gamma}$ norms bounded by a constant depending only on
$N,R,m$ and the above $C^{m-1,\gamma}$ bound for $q_{U_*}$. A finite cover of
$\partial B_r(x_{U_*})$, together with the already-known small $C^1$ slope,
turns these local estimates into a global radial graph bound
\[
\norm{g_r}_{C^{m,\gamma}(\partial B_r)}\le M_m'(N,R).
\]
The final dilation in Lemma~\ref{lem:rescale} changes this norm by at most a
bounded factor because $r\in[1/2,2]$. This gives the claimed $M_m(N,R)$ for
$\widetilde\Omega$.
\end{proof}

\begin{lemma}[Interpolation to the nearly spherical threshold]\label{lem:interp}
Fix $0<\gamma_0<\gamma$ and choose $m\ge3$. There are
$h_{\rm sph}(N,R,\gamma_0)>0$ and
\[
\alpha_{\rm sph}(N,R,\gamma_0)
:=\left(\frac{h_{\rm sph}}{C_\infty C_H'(N,R)}\right)^{2N}
\]
such that every selected graph with
$\alpha(\widetilde\Omega)\le\alpha_{\rm sph}$ satisfies
\[
\norm{g}_{C^{2,\gamma_0}(\partial B_1)}
\le \delta_{\rm sph}(N,\gamma_0),
\]
the smallness hypothesis of BDV Theorem~3.3.
\end{lemma}

\begin{proof}
Graph entry gives
\[
\norm{g}_{L^\infty(\partial B_1)}
\le C_\infty d_H(\partial\widetilde\Omega,\partial B_1)
\le C_\infty C_H'(N,R)\alpha(\widetilde\Omega)^{1/(2N)}.
\]
The interpolation inequality on $\partial B_1$ gives, for some
$\theta=\theta(N,m,\gamma_0)\in(0,1)$,
\[
\norm{g}_{C^{2,\gamma_0}}
\le C\,\norm{g}_{L^\infty}^{1-\theta}
       \norm{g}_{C^{m,\gamma}}^\theta
\le C\,\norm{g}_{L^\infty}^{1-\theta}M_m(N,R)^\theta.
\]
Choose $h_{\rm sph}$ so that the last quantity is at most
$\delta_{\rm sph}(N,\gamma_0)$ whenever $\norm{g}_{L^\infty}\le h_{\rm sph}$.
\end{proof}

\section{Closure by BDV's nearly spherical theorem}\label{sec:closure}

\begin{theorem}[BDV nearly spherical, restated as black box]\label{thm:nearlysph}
\textbf{(BDV Theorem 3.3.)} For every $0<\gamma_0\le1$ there exist
$c_{\rm sph}(N)>0$ and $\delta_{\rm sph}(N,\gamma_0)>0$ such that for every set $U$ of the form
$\partial U=\{(1+g(\theta))\theta:\theta\in\partial B_1\}$ with
\begin{itemize}
\item $\abs U=\abs{B_1}$ (volume normalization),
\item $\int_{\partial B_1}g\,\theta\d\theta=0$ (barycenter normalization),
\item $\norm{g}_{C^{2,\gamma_0}(\partial B_1)}\le\delta_{\rm sph}(N,\gamma_0)$,
\end{itemize}
we have
\[
E(U)-E(B_1)\ge c_{\rm sph}(N)\norm{g}_{H^{1/2}(\partial B_1)}^2.
\]
\end{theorem}

This is BDV Theorem 3.3 verbatim. We use it without modification.

\begin{lemma}[Asymmetry vs.\ $L^2$ graph norm, BDV Lemma 4.2 + 5.2]\label{lem:asymL2}
For sets $U$ of the form above with $\abs U=\abs{B_1}$ and
$\norm{g}_{C^0(\partial B_1)}\le\delta_0$ small,
\[
\alpha(U)\le C_{\rm BDV}(N)\norm{g}_{L^2(\partial B_1)}^2.
\]
\end{lemma}

\begin{proof}
Direct computation: $\alpha(U)=\int_U(\abs{x}-1)\d x=
\int_{\partial B_1}\int_0^{1+g}(r-1)r^{N-1}\d r\d\theta
=\int_{\partial B_1}\bigl[\tfrac12 g^2+O(g^3)\bigr]\d\theta$.
The $O(g^3)$ tail is dominated by $\norm{g}_{C^0}\norm{g}_{L^2}^2\le\delta_0\norm{g}_{L^2}^2$.
\end{proof}

\begin{corollary}[Quotient lower bound for graph minimizers]\label{cor:Qlower}
For every $U$ as in Theorem~\ref{thm:nearlysph},
\[
Q_\alpha(U)=\frac{E(U)-E(B_1)}{\alpha(U)}\ge c_*(N):=\frac{c_{\rm sph}(N)}{C_{\rm BDV}(N)}.
\]
\end{corollary}

\begin{proof}
Combine Theorem~\ref{thm:nearlysph} and Lemma~\ref{lem:asymL2}, noting
$\norm{g}_{L^2}^2\le\norm{g}_{H^{1/2}}^2$ (interpolation on $\partial B_1$).
\end{proof}

\section{The non-contradictory quantitative argument}\label{sec:argument}

We now assemble Theorem~\ref{thm:selection}, Theorem~\ref{thm:graphentry},
Proposition~\ref{prop:bootstrap}, Lemma~\ref{lem:interp}, and
Corollary~\ref{cor:Qlower} into a direct lower bound on
$Q_\alpha(\Omega_0)$.

Set
\[
\alpha_{\rm small}(N,R):=\frac12\min\{\varepsilon_{\rm sel},\alpha_{\rm graph},
\alpha_{\rm sph}\},
\]
where $\varepsilon_{\rm sel}$ is the smallness threshold in the selection
lemma and $\alpha_{\rm sph}$ is Lemma~\ref{lem:interp}'s interpolation
threshold.

\begin{theorem}[Sharp Saint--Venant stability, single-set form]\label{thm:sv-sharp}
There exists $q_*(N,R)>0$ such that for every open
$\Omega_0\subset B_R$ with $\abs{\Omega_0}=\abs{B_1}$,
$\alpha(\Omega_0)\le \alpha_{\rm small}(N,R)$, and
$Q_\alpha(\Omega_0)<q_{\rm sel}(N,R)$,
\[
Q_\alpha(\Omega_0)\ge q_*(N,R)=\frac{c_*(N)}{2C_{\rm sel}(N,R)},
\]
or equivalently
\[
E(\Omega_0)-E(B_1)\ge q_*(N,R)\,\alpha(\Omega_0).
\]
\end{theorem}

\begin{proof}
Fix $\Omega_0$ as in the hypotheses, $q_0:=Q_\alpha(\Omega_0)<q_{\rm sel}$.
Apply Theorem~\ref{thm:selection} with $\tau=q_0^{1/4}$; then
$\tau\le q_{\rm sel}^{1/4}\le\tau_{\rm reg}$ for $q_{\rm sel}$ small enough,
and $\rho(q_0,\tau)\le1/2$. We obtain a volume-normalized selected set
$\widetilde\Omega\subset B_R$ (after Lemma~\ref{lem:rescale}) with
\begin{equation}\label{eq:sel}
\alpha(\widetilde\Omega)\ge\tfrac12\alpha(\Omega_0),
\quad
\alpha(\widetilde\Omega)\le\tfrac32\alpha(\Omega_0),
\quad
Q_\alpha(\widetilde\Omega)\le 2C_{\rm sel}(N,R)\,q_0,
\end{equation}
selection parameter $\sigma\le Cq_0^{1/4}$.

\emph{Graph entry.} Since $\alpha(\Omega_0)\le\alpha_{\rm small}$,
\eqref{eq:sel} gives
\[
\alpha(\widetilde\Omega)\le \frac32\alpha_{\rm small}
\le \alpha_{\rm graph},\qquad
\alpha(\widetilde\Omega)\le\alpha_{\rm sph}.
\]
By Theorem~\ref{thm:graphentry},
$\partial\widetilde\Omega=\{(1+g)\theta\}$ is a $C^{1,\gamma}$ graph with
$\norm{g}_{C^{1,\gamma}}\le C\mu$. By Lemma~\ref{lem:interp},
$\norm{g}_{C^{2,\gamma_0}}\le\delta_{\rm sph}(N,\gamma_0)$.

\emph{Nearly spherical and quotient bound.} By
Corollary~\ref{cor:Qlower}, $Q_\alpha(\widetilde\Omega)\ge c_*(N)$.

\emph{Combining.} \eqref{eq:sel} gives
$Q_\alpha(\widetilde\Omega)\le 2C_{\rm sel}q_0$, while the previous step gives
$Q_\alpha(\widetilde\Omega)\ge c_*(N)$. Hence
\[
q_0\ge\frac{c_*(N)}{2C_{\rm sel}(N,R)}=:q_*(N,R).
\]
\end{proof}

\begin{remark}[Why this is not a contradiction]
At no point did we assume $Q_\alpha(\Omega_0)$ is arbitrarily small. We
\emph{computed} the inequality
$Q_\alpha(\Omega_0)\ge c_*(N)/(2C_{\rm sel}(N,R))$
directly from the chain of constants. The selection is single-set; the
graph entry is below an explicit threshold $\alpha_{\rm graph}(N,R)$; the
closure is the pointwise BDV nearly spherical inequality.

The compactness/limit step of BDV is replaced by Theorem~\ref{thm:selection},
which produces $\widetilde\Omega$ from a single $\Omega_0$ with explicit
$(N,R)$-constants.
\end{remark}

\subsection{Extension to all asymmetries}

\begin{theorem}[Sharp Saint--Venant, global form]\label{thm:sv-global}
There exists $c_{\rm SV}(N,R)>0$ such that for every open
$\Omega\subset B_R$ with $\abs{\Omega}=\abs{B_1}$,
\[
E(\Omega)-E(B_1)\ge c_{\rm SV}(N,R)\,\mathcal A(\Omega)^2.
\]
Explicitly,
\[
c_{\rm SV}(N,R)
=\min\Bigl(\frac{q_*(N,R)}{C_{\rm Fr}(N)},\;c_{\rm far,SV}(N,R)\Bigr),
\]
where $C_{\rm Fr}(N)$ is BDV Lemma 4.2's constant and $c_{\rm far,SV}(N,R)$
is the explicit far-from-ball lower bound from \S\ref{sec:far} below.
\end{theorem}

\begin{proof}
Small asymmetry, $\alpha(\Omega)\le\alpha_{\rm small}$:
combine Theorem~\ref{thm:sv-sharp} with $\mathcal A^2\le C_{\rm Fr}\alpha$.

Large asymmetry, $\alpha(\Omega)>\alpha_{\rm small}$:
direct compactness, see \S\ref{sec:far}.

Take $c_{\rm SV}:=\min$ of the two constants.
\end{proof}

\subsection{The far-from-ball Saint--Venant constant}\label{sec:far}

\begin{lemma}[Far-from-ball uniform lower bound]\label{lem:far}
There exists $c_{\rm far,SV}(N,R)>0$ such that for every open
$\Omega\subset B_R$ with $\abs{\Omega}=\abs{B_1}$ and
$\alpha(\Omega)\ge\alpha_{\rm small}(N,R)$,
\[
E(\Omega)-E(B_1)\ge c_{\rm far,SV}(N,R).
\]
\end{lemma}

\begin{proof}
First convert $\alpha$-separation into Fraenkel separation. Let $B$ be an
optimal ball for $\mathcal A(\Omega)$. Since $\Omega\subset B_R$ and
$|B|=|B_1|$, the relevant balls are contained in a fixed ball $B_{R+2}$ after
enlarging $R$ by a dimensional amount. BDV Lemma~4.2(ii), applied to
$\Omega$ and $B$, gives
\[
\alpha(\Omega)\le C_2(N,R)\,|\Omega\Delta B|
=C_2(N,R)|B_1|\mathcal A(\Omega).
\]
Thus $\alpha(\Omega)\ge\alpha_{\rm small}$ implies
\[
\mathcal A(\Omega)\ge
\alpha_{\rm small}/(C_2(N,R)|B_1|)=:a_{\rm far}(N,R)>0.
\]
Now use the suboptimal quantitative Saint--Venant inequality quoted by BDV as
Lemma~5.2:
\[
E(\Omega)-E(B_1)\ge C_8(N)^{-1}\mathcal A(\Omega)^4.
\]
Therefore one can take
\[
c_{\rm far,SV}(N,R):=C_8(N)^{-1}a_{\rm far}(N,R)^4.
\]
\end{proof}

\begin{remark}
The Lemma~\ref{lem:far} infimum depends on $(N,R)$ but is not given
explicitly. For numerical applications, $c_{\rm far,SV}$ can be estimated by
known sharp inequalities for specific classes (convex bodies, smooth
domains), but for the theorem above we only need positivity.
\end{remark}

\section{Kohler--Jobin transfer to Faber--Krahn}\label{sec:kj}

\begin{theorem}[Kohler--Jobin inequality, 1978]\label{thm:kj}
For every open bounded $\Omega\subset\R^N$ with $\lambda_1(\Omega)<\infty$,
\[
\Psi(\Omega):=T(\Omega)\cdot\lambda_1(\Omega)^{(N+2)/2}\ge\Psi(B),
\]
where $B$ is any ball, with equality iff $\Omega$ is a ball.
\end{theorem}

\begin{lemma}[Pointwise transfer at fixed volume]\label{lem:transferpt}
For every open $\Omega$ with $\abs{\Omega}=\abs{B^*}$,
\[
\lambda_1(\Omega)\ge\lambda_1(B^*)\Bigl(\frac{T(B^*)}{T(\Omega)}\Bigr)^{2/(N+2)}.
\]
\end{lemma}

\begin{lemma}[Small-deficit linearization]\label{lem:transferlin}
Let $\delta_{\rm SV}:=T(B^*)-T(\Omega)$. For $\delta_{\rm SV}\le T(B^*)/2$,
\[
\lambda_1(\Omega)-\lambda_1(B^*)\ge\frac{2\lambda_1(B^*)}{(N+2)T(B^*)}\,\delta_{\rm SV}.
\]
\end{lemma}

\begin{proof}
For $t\in[0,1/2]$, $(1-t)^{-1/\beta}-1\ge t/\beta$ where $\beta=(N+2)/2$.
Apply with $t=\delta_{\rm SV}/T(B^*)$ and use $\lambda_1(\Omega)\ge L\cdot(1-t)^{-1/\beta}$.
\end{proof}

\begin{remark}[The factor $2L_N/((N+2)T_N)$]
$L_N:=\lambda_1(B_1)=j_{N/2-1,1}^2$ where $j_{\nu,1}$ is the first positive
zero of the Bessel function $J_\nu$. $T_N:=T(B_1)=\abs{B_1}/(N(N+2))$.
For $N=2$: $L_2=j_{0,1}^2\approx5.78$, $T_2=\pi/8$, multiplier
$\approx2\cdot5.78/(4\cdot\pi/8)\approx3.68/\pi$. For $N=3$: $L_3=\pi^2$,
$T_3=4\pi/15$, multiplier $\approx2\pi^2/(5\cdot 4\pi/15)=3\pi/2$.
These are universal in $N$, with no $R$-dependence.
\end{remark}

\section{The sharp Faber--Krahn theorem}\label{sec:final}

\begin{theorem}[Sharp Faber--Krahn stability, Route A end-to-end]\label{thm:fk-final}
There exists $c_{\rm FK}(N,R)>0$ such that for every open
$\Omega\subset B_R$ with $\abs{\Omega}=\abs{B^*}$,
\[
\lambda_1(\Omega)-\lambda_1(B^*)\ge c_{\rm FK}(N,R)\,\mathcal A(\Omega)^2.
\]
\end{theorem}

\begin{proof}
\emph{Small asymmetry.} By Theorem~\ref{thm:sv-global}, in the small-asymmetry
regime $\mathcal A(\Omega)\le\mathcal A_*(N,R)$,
$\delta_{\rm SV}:=T(B^*)-T(\Omega)=2(E(\Omega)-E(B^*))
\ge 2c_{\rm SV}(N,R)\mathcal A(\Omega)^2$
(where $E=-T/2$). By Lemma~\ref{lem:transferlin}, for $\delta_{\rm SV}$ small,
\[
\lambda_1(\Omega)-\lambda_1(B^*)\ge\frac{2L_N}{(N+2)T_N}\,\delta_{\rm SV}
\ge\frac{4L_N\,c_{\rm SV}(N,R)}{(N+2)T_N}\,\mathcal A(\Omega)^2.
\]

\emph{Large asymmetry or large torsion deficit.} If
$\mathcal A(\Omega)>\mathcal A_*$, the suboptimal Faber--Krahn stability
quoted by BDV gives a positive lower bound depending only on $(N,R)$; call it
$c_{\rm far,FK}(N,R)$. If instead $\delta_{\rm SV}>T_N/2$, the Kohler--Jobin
pointwise bound gives
\[
\lambda_1(\Omega)-L_N\ge L_N(2^{2/(N+2)}-1),
\]
which is also bounded below by a dimensional multiple of
$\mathcal A(\Omega)^2$ because $\mathcal A\le2$.

Take
\[
c_{\rm FK}(N,R):=\min\Bigl(\frac{4L_N\,c_{\rm SV}(N,R)}{(N+2)T_N},\;\frac{c_{\rm far,FK}(N,R)}{4},
\frac{L_N(2^{2/(N+2)}-1)}{4}\Bigr).
\]
\end{proof}

\begin{remark}[Full constant chain]
The constants are produced in the order:
\begin{enumerate}
\item Selection: $C_{\rm sel}(N,R), q_{\rm sel}(N,R), \tau_{\rm reg}(N,R)$ from
\texttt{quantitative-selection-principle.md} (depending on BDV Lemmas 4.9, 4.12, 4.16).
\item Graph entry: $\alpha_{\rm graph}(N,R)$ from
$\alpha_{\rm graph}=(h_F/C_H')^{2N}$,
where $h_F=\mu\rho_\mu/16$ and $\mu,\rho_\mu$ are BDV Theorem~4.18 thresholds.
\item Bootstrap: $C_{\rm boot}(N,R)$ from Kinderlehrer--Nirenberg /
Lieberman applied with BDV Lemma 4.16's $L_q$.
\item Nearly spherical: $c_{\rm sph}(N), \delta_{\rm sph}(N,\gamma_0)$ from BDV
Theorem~3.3.
\item Asymmetry vs.\ $L^2$: $C_{\rm BDV}(N)$ from BDV Lemma~4.2 + 5.2.
\item Saint--Venant quotient: $c_*(N)=c_{\rm sph}(N)/C_{\rm BDV}(N)$.
\item Saint--Venant constant: $q_*(N,R)=c_*(N)/(2C_{\rm sel}(N,R))$.
\item Global Saint--Venant: $c_{\rm SV}(N,R)=\min(q_*/C_{\rm Fr},c_{\rm far,SV})$.
\item Kohler--Jobin multiplier: $2L_N/((N+2)T_N)$, universal.
\item Faber--Krahn: $c_{\rm FK}(N,R)$ is the minimum displayed in
Theorem~\ref{thm:fk-final}.
\end{enumerate}
Every constant in this chain depends only on $(N,R)$, and the dependence is
through named BDV regularity constants. No constant arises from a limit or
compactness argument that hides its $(N,R)$-dependence; the only
far-regime constants are $c_{\rm far,SV}$ and $c_{\rm far,FK}$, which are
obtained from suboptimal stability bounds away from the ball and do not enter
the small-asymmetry regime.
\end{remark}

\section{Where the contradiction was, and what replaced it}

\begin{center}
\begin{tabular}{lll}
\hline
Step & BDV & Plan 1 / Route A \\
\hline
Selection & sequential + compactness limit (contradiction) &
  single-set + explicit $C_{\rm sel}(N,R)$ \\
Limit existence & qualitative, $U_n\to U_\infty$ in $L^1$ &
  not needed \\
$\alpha\to$ graph & qualitative from limit & explicit $\alpha_{\rm graph}(N,R)$ \\
Bootstrap & implicit in BDV regularity & explicit Schauder
  (Kinderlehrer--Nirenberg) \\
Nearly spherical & BDV Theorem 3.3 & unchanged \\
Quotient lower bound & contradiction with $\alpha(U_\infty)>0$ &
  direct, $Q_\alpha\ge c_*(N)$ \\
Saint--Venant & by contradiction & by computation,
  $Q_\alpha\ge c_*/2C_{\rm sel}$ \\
Faber--Krahn transfer & Kohler--Jobin & unchanged \\
\hline
\end{tabular}
\end{center}

\paragraph{The single contradiction step that was removed.}
BDV's argument is:
``Assume sharp Saint--Venant fails $\Rightarrow$ minimizing sequence
$\Rightarrow$ compactness limit $\Rightarrow$ contradiction with nearly spherical.''

Plan 1's argument is:
``Given any $\Omega_0$ with $\alpha(\Omega_0)\le\alpha_{\rm small}$:
\begin{itemize}
\item Construct $\widetilde\Omega$ in one step (Thm.~\ref{thm:selection}).
\item Note $\widetilde\Omega$ enters the small $C^{2,\gamma_0}$ nearly
spherical class (Thm.~\ref{thm:graphentry}, Prop.~\ref{prop:bootstrap},
Lemma~\ref{lem:interp}).
\item Apply BDV Theorem 3.3 pointwise to read off
$Q_\alpha(\widetilde\Omega)\ge c_*(N)$.
\item Conclude $Q_\alpha(\Omega_0)\ge c_*(N)/(2C_{\rm sel}(N,R))$
by the selection quotient bound.''
\end{itemize}
No assumption on smallness of $Q_\alpha(\Omega_0)$ is made; the conclusion
is a direct lower bound.

\section{Status}

\paragraph{What is fully written.}
\begin{itemize}
\item Theorem~\ref{thm:selection} (selection): full proof in
\texttt{quantitative-selection-principle.md}.
\item Lemma~\ref{lem:rescale} (volume normalization): full proof in
\texttt{corrections-response.tex} \S 1.
\item Lemmas~\ref{lem:alphaH}, Props.~\ref{prop:flatness},
\ref{prop:globalgraph}, Thm.~\ref{thm:graphentry} (graph entry):
full proofs in \texttt{graph-entry-closure.md}.
\item Prop.~\ref{prop:bootstrap} and Lemma~\ref{lem:interp} (uniform
Schauder bounds plus interpolation to small $C^{2,\gamma_0}$).
\item Theorem~\ref{thm:nearlysph}: BDV Theorem 3.3, used as black box.
\item Lemmas~\ref{lem:asymL2}, \ref{lem:transferpt}, \ref{lem:transferlin}:
elementary, written in this note.
\item Theorems~\ref{thm:sv-sharp}, \ref{thm:sv-global}, \ref{thm:fk-final}:
assembled in this note.
\end{itemize}

\paragraph{What uses BDV as a black box.}
BDV Lemmas 4.9 (density), 4.12 (Lipschitz/nondegeneracy), 4.16 (Bernoulli
coefficient), 4.18 (Alt--Caffarelli flatness improvement), 3.3 (nearly
spherical), 4.2/5.2 ($\alpha$ vs $L^2$ and suboptimal stability). All cited by theorem number and used
without modification.

\paragraph{What is not numerically tracked.}
The explicit values of $C_{\rm sel}, q_{\rm sel}, \alpha_{\rm graph},
C_{\rm boot}, c_{\rm sph}, C_{\rm BDV}, c_{\rm far,SV}, c_{\rm far,FK}$
are not computed. The point of Route A is the \emph{structural}
quantification: every constant is named and $(N,R)$-trackable, but
extracting concrete numbers requires going through each BDV proof.

\paragraph{What this delivers over BDV.}
\begin{enumerate}
\item BDV's contradiction step (sequential selection + compactness limit)
is replaced by Theorem~\ref{thm:selection}, a one-step construction with
explicit constants.
\item The graph-entry threshold $\alpha_{\rm graph}(N,R)$ is given by a
formula.
\item The constant $c_{\rm FK}(N,R)$ is the minimum of a chain of named
quantities, traceable to the BDV regularity package.
\item The proof is amenable to numerical/effective tracking: one can take
each BDV proof and extract explicit constants, then plug through this
chain.
\end{enumerate}

\paragraph{Auditability.}
Every lemma in this note is either:
\begin{itemize}
\item proved in a referenced file with the same notation,
\item proved elementarily in this note (Lemmas~\ref{lem:asymL2},
\ref{lem:transferpt}, \ref{lem:transferlin}, far-from-ball
Lemma~\ref{lem:far}), or
\item quoted from BDV by theorem number (Thm.~3.3, Lemmas~4.2, 4.9, 4.12,
4.16, 5.2, and Thm.~4.18).
\end{itemize}
There are no appeals to ``a standard argument'' that is not pinned to a
reference.

\end{document}
